\documentclass[12pt,a4paper]{article}

% Packages
\usepackage[utf8]{inputenc}
\usepackage[T1]{fontenc}
\usepackage[french]{babel}
\usepackage{geometry}
\usepackage{hyperref}
\usepackage{graphicx}
\usepackage{listings}
\usepackage{xcolor}
\usepackage{titlesec}
\usepackage{fancyhdr}

% Geometry
\geometry{
    left=2.5cm,
    right=2.5cm,
    top=3cm,
    bottom=3cm,
    headheight=15pt
}

% Listings configuration for code
\lstset{
    basicstyle=\ttfamily\small,
    breaklines=true,
    frame=single,
    numbers=left,
    numberstyle=\tiny\color{gray},
    keywordstyle=\color{blue},
    commentstyle=\color{green!60!black},
    stringstyle=\color{orange},
    showstringspaces=false,
    tabsize=4
}

% YAML language definition for listings
\lstdefinelanguage{yaml}{
    keywords={true,false,null,yes,no,on,off},
    keywordstyle=\color{blue}\bfseries,
    keywords=[2]{stage,script,image,tags,before_script,after_script,artifacts,
                 cache,rules,needs,variables,when,retry,default,stages,
                 include,extends,only,except,allow_failure,timeout,coverage,
                 services,dependencies,trigger,parallel,resource_group,
                 environment,deployment,pages,interruptible,retry,release},
    keywordstyle=[2]\color{purple}\bfseries,
    sensitive=true,
    comment=[l]{\#},
    commentstyle=\color{green!60!black}\itshape,
    morestring=[b]',
    morestring=[b]",
    stringstyle=\color{orange},
    literate=
        {-\ }{{\textcolor{blue}{-\ }}}{2}
        {>}{{\textcolor{red}{>}}}{1}
        {|}{{\textcolor{red}{|}}}{1}
        {:\ }{{\textcolor{red}{:\ }}}{2}
        {\ -\ }{{\textcolor{blue}{\ -\ }}}{3}
        {\$\{}{{\textcolor{teal}{\$\{}}}{2}
        {\}}{{\textcolor{teal}{\}}}}{1}
        {\$}{{\textcolor{teal}{\$}}}{1}
        {||}{{\textcolor{purple}{||}}}{2}
        {\&\&}{{\textcolor{purple}{\&\&}}}{2}
        {==}{{\textcolor{purple}{==}}}{2}
        {!=}{{\textcolor{purple}{!=}}}{2},
    basicstyle=\ttfamily\small\color{black},
    breaklines=true,
    showstringspaces=false,
    tabsize=2
}

% Enhanced XML language definition for listings
\lstdefinelanguage{enhancedxml}{
    basicstyle=\ttfamily\small\color{black},
    morestring=[b]",
    morestring=[s]{>}{<},
    morecomment=[s]{<?}{?>},
    morecomment=[s]{<!--}{-->},
    commentstyle=\color{green!60!black}\itshape,
    stringstyle=\color{orange},
    identifierstyle=\color{blue},
    keywordstyle=\color{purple}\bfseries,
    morekeywords={xmlns,version,type,encoding,id,url,name,value,phase,goal,
                  configuration,execution,executions,plugin,plugins,build,
                  project,dependencies,dependency,groupId,artifactId,
                  version,scope,properties,property,repository,repositories,
                  distributionManagement,snapshotRepository,server,servers,
                  username,password,tagNameFormat,autoVersionSubmodules,
                  releaseProfiles,scmCommentPrefix,goals},
    sensitive=true,
    alsoletter={:,-},
    literate=
        {<}{{\textcolor{red}{<}}}{1}
        {>}{{\textcolor{red}{>}}}{1}
        {</}{{\textcolor{red}{</}}}{2}
        {/>}{{\textcolor{red}{/>}}}{2}
        {<?}{{\textcolor{purple}{<?}}}{2}
        {?>}{{\textcolor{purple}{?>}}}{2}
        {\$\{env.}{{\textcolor{teal}{\$\{env.}}}{6}
        {\$\{project.}{{\textcolor{teal}{\$\{project.}}}{10}
        {\}}{{\textcolor{teal}{\}}}}{1},
    breaklines=true,
    showstringspaces=false,
    tabsize=4
}

% Header and footer
\pagestyle{fancy}
\fancyhf{}
\rhead{Rapport de Projet - Groupe 7}
\lhead{MiniJAJA}
\rfoot{Page \thepage}

% Title information
\title{
    \textbf{Rapport de Projet}\\
    \large{MiniJAJA - Compilateur et Interpréteur}\\
    \large{Groupe 7}
}
\author{
    Romain \and
    Lucas \and
    Félix \and
    Léo \and
    Théo \and
    Ahmed
}
\date{\today}

\begin{document}

\maketitle
\newpage

\tableofcontents
\newpage


% ============================================================================
% SECTION : Traitement des erreurs
% ============================================================================
\section{Traitement des erreurs}\label{sec:traitement-des-erreurs}

\subsection{Introduction}\label{subsec:introduction}
Cette section présente les principaux choix de réalisation concernant le traitement des erreurs dans le projet MiniJAJA. Le traitement des erreurs est un aspect crucial du compilateur et de l'interpréteur, permettant de fournir des retours clairs et précis aux utilisateurs lors de la détection d'anomalies.

\subsection{Stratégie générale de gestion des erreurs}\label{subsec:strategie-generale-de-gestion-des-erreurs}

\subsubsection{Classification des erreurs}
Le projet MiniJAJA gère plusieurs types d'erreurs à travers un système de diagnostic unifié.
Le système distingue quatre phases d'erreurs, définies dans l'énumération \texttt{Phase} :
\begin{itemize}
    \item \textbf{Erreurs lexicales} (\texttt{LEXICAL}) : détectées lors de l'analyse lexicale (tokens invalides, caractères non reconnus)
    \item \textbf{Erreurs syntaxiques} (\texttt{SYNTAX}) : détectées lors de l'analyse syntaxique (règles de grammaire violées)
    \item \textbf{Erreurs sémantiques} (\texttt{SEMANTIC}) : détectées lors de l'analyse sémantique (typage, portée, déclarations)
    \item \textbf{Erreurs d'exécution} (\texttt{RUNTIME}) : détectées lors de l'interprétation du code JajaCode
\end{itemize}

Chaque erreur possède également une sévérité définie par l'énumération \texttt{Severity} :
\begin{itemize}
    \item \textbf{ERROR} : erreur bloquante empêchant la compilation ou l'exécution
    \item \textbf{WARNING} : avertissement n'empêchant pas la compilation, mais signalant un problème potentiel.
    \textit{Note : cette fonctionnalité est définie dans l'architecture, mais n'est pas encore exploitée dans l'implémentation actuelle. L'infrastructure est en place pour permettre une utilisation future.}
\end{itemize}

\subsubsection{Architecture du système de diagnostic}

Le système de gestion des erreurs repose sur une architecture modulaire composée de plusieurs classes clés :

\paragraph{Classe Diagnostic}
La classe \texttt{Diagnostic} est un record Java qui encapsule toutes les informations d'une erreur :

\vspace{0.3cm}
\begin{lstlisting}[language=java, caption=Structure du Diagnostic,label={lst:diagnostic-record}]
public record Diagnostic(
    Severity severity,      // ERROR ou WARNING
    Phase phase,            // LEXICAL, SYNTAX, SEMANTIC ou RUNTIME
    SourcePosition position, // Fichier, ligne, colonne
    String message          // Message descriptif
) {}
\end{lstlisting}

\paragraph{Classe SourcePosition}
La classe \texttt{SourcePosition} capture la position exacte de l'erreur dans le code source, incluant le nom du fichier, le numéro de ligne et la position dans la ligne.
Cette information permet aux développeurs de localiser rapidement l'origine du problème.

\paragraph{Classe DiagnosticCollector}
Le \texttt{DiagnosticCollector} est responsable de collecter tous les diagnostics durant les phases d'analyse.
Il permet de :
\begin{itemize}
    \item Accumuler plusieurs erreurs sans interrompre l'analyse
    \item Formater tous les diagnostics en une seule chaîne de caractères via \texttt{formatDiagnostics()}
    \item Vérifier la présence d'erreurs spécifiques par phase :
    \begin{itemize}
        \item \texttt{hasErrors()} : présence d'erreurs de type ERROR
        \item \texttt{hasSyntaxErrors()} : erreurs syntaxiques ou lexicales (phases SYNTAX ou LEXICAL)
        \item \texttt{hasSemanticErrors()} : erreurs sémantiques (phase SEMANTIC)
    \end{itemize}
\end{itemize}

\subsubsection{Principes de conception}
Les choix de conception pour le traitement des erreurs suivent plusieurs principes :
\begin{itemize}
    \item \textbf{Clarté} : les messages d'erreur sont explicites et compréhensibles
    \item \textbf{Précision} : indication exacte de la position (fichier, ligne, colonne)
    \item \textbf{Récupération} : poursuite de l'analyse pour détecter plusieurs erreurs simultanément
    \item \textbf{Cohérence} : uniformité du format des messages via le système \texttt{Diagnostic}
    \item \textbf{Traçabilité} : chaque erreur est associée à sa phase de détection
\end{itemize}

\subsection{Traitement des erreurs par module}\label{subsec:traitement-des-erreurs-par-module}

\subsubsection{Module LexerParser}

\paragraph{Erreurs lexicales et syntaxiques}

Le module LexerParser utilise ANTLR4 pour l'analyse lexicale et syntaxique.
Les erreurs sont capturées via un listener personnalisé :

\begin{lstlisting}[language=java, caption=SyntaxErrorListener,label={lst:syntaxerrorlistener}]
public class SyntaxErrorListener extends BaseErrorListener {
    private final DiagnosticCollector diagnosticCollector;
    private final String fileName;

    @Override
    public void syntaxError(Recognizer<?, ?> recognizer,
                           Object offendingSymbol,
                           int line, int charPositionInLine,
                           String msg, RecognitionException e) {
        SourcePosition pos = new SourcePosition(fileName,
                                                line,
                                                charPositionInLine);
        diagnosticCollector.report(Severity.ERROR,
                                  Phase.SYNTAX,
                                  pos, msg);
    }
}
\end{lstlisting}

Le \texttt{SyntaxErrorListener} :
\begin{itemize}
    \item Remplace les listeners d'erreur par défaut d'ANTLR
    \item Capture les erreurs lexicales et syntaxiques
    \item Les transforme en objets \texttt{Diagnostic} structurés
    \item Permet la collecte de multiples erreurs sans interruption
\end{itemize}

\paragraph{Erreurs sémantiques}

L'analyse sémantique est réalisée par le \texttt{MiniJajaSemanticAnalyzer}, qui délègue les vérifications à des composants spécialisés :

\begin{itemize}
    \item \textbf{TypeChecker} : vérifie la compatibilité des types dans les expressions et affectations
    \item \textbf{ScopeResolver} : résout les références aux variables et vérifie leur portée
    \item \textbf{DeclarationCollector} : collecte les déclarations et détecte les redéclarations
    \item \textbf{MethodCallValidator} : valide les appels de méthodes (existence, arité, types des paramètres)
    \item \textbf{TypeInferenceEngine} : infère les types des expressions complexes
\end{itemize}

Les erreurs sémantiques typiques incluent :
\begin{itemize}
    \item Variables non déclarées
    \item Redéclaration de variables
    \item Incompatibilités de types (ex: affectation d'un entier à un booléen)
    \item Appels de méthodes inexistantes
    \item Mauvais nombre d'arguments dans les appels de méthode
    \item Accès à des tableaux avec un indice non entier
    \item Utilisation de variables hors de leur portée
\end{itemize}

Chaque validateur utilise un \texttt{SemanticContext} partagé qui maintient :
\begin{itemize}
    \item La table des symboles (variables, méthodes, constantes)
    \item Les informations de typage
    \item La portée courante
\end{itemize}

\subsubsection{Interpréteur JajaCode}

L'interpréteur JajaCode utilise une hiérarchie d'exceptions spécifiques pour les erreurs d'exécution, toutes héritant de \texttt{JajaCodeRuntimeException}.

\paragraph{Exception de base : JajaCodeRuntimeException}

\vspace{0.3cm}
\begin{lstlisting}[language=java, caption=JajaCodeRuntimeException,label={lst:jjcruntimeexception}]
public class JajaCodeRuntimeException extends RuntimeException {
    private final int programCounter;
    private final String axiomeName;

    public JajaCodeRuntimeException(String message,
                                    String axiomeName,
                                    int programCounter) {
        super(String.format("[PC=%d, Axiome=%s] %s",
                           programCounter, axiomeName, message));
        this.axiomeName = axiomeName;
        this.programCounter = programCounter;
    }
}
\end{lstlisting}

Cette exception de base fournit :
\begin{itemize}
    \item Le \textbf{Program Counter (PC)} : position dans le code JajaCode où l'erreur s'est produite.
    \item Le \textbf{nom de l'axiome} : l'instruction JajaCode en cours d'exécution.
    \item Un \textbf{message formaté} incluant ces informations contextuelles.
\end{itemize}

\paragraph{Exceptions spécifiques}

Le projet définit plusieurs exceptions spécialisées :

\begin{itemize}
    \item \textbf{DivisionByZeroException} : division ou modulo par zéro lors d'opérations arithmétiques,
    \item \textbf{StackUnderflowException} : tentative de pop sur une pile vide,
    \item \textbf{TypeMismatchException} : incompatibilité de types à l'exécution,
    \item \textbf{InvalidAddressException} : accès mémoire à une adresse invalide,
    \item \textbf{AssignmentException} : erreur lors d'une affectation,
    \item \textbf{UndefinedSymbolException} : référence à un symbole non défini.
\end{itemize}

Exemple d'exception spécialisée :
\begin{lstlisting}[language=java, caption=DivisionByZeroException,label={lst:divisionbyzero-exception}]
public class DivisionByZeroException
    extends JajaCodeRuntimeException {

    public DivisionByZeroException(String axiomeName,
                                   int programCounter) {
        super("Division par zero", axiomeName, programCounter);
    }
}
\end{lstlisting}

\subsection{Format des messages d'erreur}\label{subsec:format-des-messages-d'erreur}

\subsubsection{Structure des messages}

Les messages d'erreur suivent un format uniforme qui dépend de la phase :

\paragraph{Erreurs de compilation (Lexicales, Syntaxiques, Sémantiques)}

\vspace{0.3cm}
\begin{lstlisting}[caption=Format des erreurs de compilation,label={lst:ex-err-compil}]
[PHASE] fichier.mjj:ligne:colonne - Message descriptif
\end{lstlisting}

Exemple :
\begin{lstlisting}[caption=Exemple d'erreur sémantique,label={lst:ex-err-semantique}]
[SEMANTIC] test.mjj:15:8 - Variable 'x' non declaree
\end{lstlisting}

\paragraph{Erreurs d'exécution (Runtime JajaCode)}

\vspace{0.3cm}
\begin{lstlisting}[caption=Format des erreurs d'execution,label={lst:ex-err-runtime}]
[PC=<compteur>, Axiome=<nom>] Message descriptif
\end{lstlisting}

Exemple :
\begin{lstlisting}[caption=Exemple d'erreur d'execution,label={lst:ex-err-execution}]
[PC=42, Axiome=div] Division par zero
\end{lstlisting}

\subsubsection{Informations contextuelles}

Chaque message d'erreur inclut des informations contextuelles permettant au développeur de localiser et comprendre le problème :
\begin{itemize}
    \item \textbf{Nom du fichier} : pour les projets multi-fichiers
    \item \textbf{Ligne et colonne} : position exacte dans le source
    \item \textbf{Program Counter} : pour les erreurs JajaCode, position dans le bytecode
    \item \textbf{Nom de l'axiome} : instruction JajaCode en cours d'exécution
    \item \textbf{Message descriptif} : explication claire du problème
\end{itemize}

\subsection{Gestion des exceptions}\label{subsec:gestion-des-exceptions}

\subsubsection{Hiérarchie des exceptions}

Le projet utilise une hiérarchie d'exceptions bien structurée organisée en trois catégories principales :

\begin{lstlisting}[caption=Hierarchie des exceptions,label={lst:hierarchie-des-exceptions}]
java.lang.RuntimeException
  |
  +-- SyntaxException (erreurs syntaxiques et lexicales)
  |
  +-- SemanticException (erreurs semantiques)
  |
  +-- JajaCodeRuntimeException (erreurs d'execution)
        |
        +-- DivisionByZeroException
        +-- StackUnderflowException
        +-- TypeMismatchException
        +-- InvalidAddressException
        +-- AssignmentException
        +-- UndefinedSymbolException
\end{lstlisting}

Cette hiérarchie permet :
\begin{itemize}
    \item Une capture sélective des exceptions par phase de compilation
    \item Un traitement différencié selon le type d'erreur
    \item Une extension facile avec de nouvelles exceptions
    \item Une séparation claire entre erreurs de compilation et d'exécution
\end{itemize}

\paragraph{Exceptions de compilation}

Les exceptions \texttt{SyntaxException} et \texttt{SemanticException} sont définies dans le package \texttt{fr.ufrst.m1info.gl.groupe7.lexerparser.errors.exceptions} et suivent un pattern unifié :

\begin{lstlisting}[language=java, caption=Pattern des exceptions de compilation,label={lst:compilation-exception-pattern}]
public class SemanticException extends RuntimeException {
    public SemanticException(DiagnosticCollector collector) {
        super(collector != null && collector.hasErrors()
                ? collector.formatDiagnostics()
                : "Unknown semantic error.");
    }
}
\end{lstlisting}

Ce pattern présente plusieurs avantages :
\begin{itemize}
    \item \textbf{Intégration avec DiagnosticCollector} : le message d'exception contient automatiquement tous les diagnostics formatés
    \item \textbf{Messages riches} : inclusion de toutes les erreurs détectées, pas seulement la première
    \item \textbf{Cohérence} : format uniforme pour toutes les exceptions de compilation
    \item \textbf{Robustesse} : message par défaut si le collector est null ou vide
\end{itemize}

\texttt{SyntaxException} est levée lorsque des erreurs lexicales ou syntaxiques sont détectées (phase LEXICAL ou SYNTAX), tandis que \texttt{SemanticException} est levée pour les erreurs d'analyse sémantique (phase SEMANTIC).

\subsubsection{Propagation et capture}

Les exceptions sont propagées jusqu'au point d'entrée principal où elles sont capturées et transformées en messages utilisateur. Le système de \texttt{DiagnosticCollector} permet également de collecter plusieurs erreurs avant de les rapporter toutes ensemble.

\subsection{Tests et validation}\label{subsec:tests-et-validation}

\subsubsection{Couverture des tests}

Le module LexerParser contient \textbf{38 classes de tests}, dont plusieurs dédiées spécifiquement au traitement des erreurs :

\begin{itemize}
    \item \texttt{ExceptionTypeTest} : vérifie que les bonnes exceptions sont levées selon le type d'erreur (syntaxique vs sémantique)
    \item \texttt{SyntaxErrorListenerTest} : teste la capture des erreurs syntaxiques via le listener ANTLR
    \item \texttt{MiniJajaSemanticAnalyzerTest} : teste l'analyse sémantique générale
    \item \texttt{MiniJajaSemanticAnalyzerScopeTest} : teste la vérification de portée des variables
    \item \texttt{MiniJajaSemanticAnalyzerReturnTest} : teste la validation des instructions de retour
    \item \texttt{DiagnosticCollectorTest} : teste le collecteur de diagnostics et ses méthodes de filtrage
    \item \texttt{DiagnosticTest} : teste la structure des diagnostics (record)
    \item \texttt{JajaCodeExceptionsTest} : teste les exceptions d'exécution JajaCode
    \item \texttt{ArrayAssignmentTest} : teste les erreurs d'affectation de tableaux
    \item \texttt{TypeIncompatibilityTest} : teste les incompatibilités de types
\end{itemize}

\subsubsection{Stratégie de test}

Les tests vérifient :
\begin{itemize}
    \item La \textbf{détection correcte} des erreurs (pas de faux négatifs)
    \item Le \textbf{type d'exception} levé selon la phase d'erreur
    \item Le \textbf{format des messages} d'erreur intégrés aux exceptions
    \item La \textbf{position exacte} (ligne/colonne) des erreurs
    \item La \textbf{phase correcte} de détection
    \item La \textbf{récupération} après erreur (continuation de l'analyse)
\end{itemize}

Exemple de test vérifiant le type d'exception :
\begin{lstlisting}[language=java, caption=Test de type d'exception semantique,label={lst:test-semantic-exception}]
@Test
@DisplayName("Semantic error should throw SemanticException")
void semanticError_shouldThrowSemanticException() {
    String codeWithSemanticError = """
        class C {
            main {
                int x = 5;
                int x = 10;  // Redeclaration
            }
        }
    """;

    assertThrows(SemanticException.class,
        () -> MiniJajaCompiler
            .getMiniJajaCompilerVisitorFromString(codeWithSemanticError),
        "Semantic errors should throw SemanticException");
}
\end{lstlisting}

Ce test vérifie que la redéclaration d'une variable lève bien une \texttt{SemanticException} et non une \texttt{SyntaxException}. D'autres tests similaires vérifient les incompatibilités de types, les variables non déclarées, et les erreurs syntaxiques.

\subsection{Améliorations futures}\label{subsec:ameliorations-futures}

Plusieurs améliorations sont envisageables pour le système de gestion des erreurs :

\begin{itemize}
    \item \textbf{Activation des warnings} : exploiter le niveau de sévérité \texttt{WARNING} déjà présent dans l'architecture pour signaler des problèmes potentiels sans bloquer la compilation. Exemples d'usage :
    \begin{itemize}
        \item Variables déclarées mais jamais utilisées
        \item Code mort (unreachable code)
        \item Comparaisons redondantes
        \item Conventions de nommage non respectées
        \item Déprécations
    \end{itemize}
    \item \textbf{Suggestions de correction} : proposer des corrections automatiques (ex: ``Vouliez-vous dire 'integer' ?'')
    \item \textbf{Messages multilingues} : support de l'internationalisation des messages
    \item \textbf{Colorisation} : messages avec codes couleur dans la CLI pour une meilleure lisibilité
    \item \textbf{Extraits de code} : afficher l'extrait de code source avec l'erreur soulignée
    \item \textbf{Erreurs catégorisées} : codes d'erreur numérotés pour une documentation détaillée
    \item \textbf{Mode strict/permissif} : niveau de sévérité configurable (warnings as errors, ignorer certains warnings)
\end{itemize}

\subsection{Conclusion}\label{subsec:conclusion}

Le système de gestion des erreurs du projet MiniJAJA repose sur une architecture robuste et extensible.
L'utilisation d'un système de diagnostic unifié à travers toutes les phases (lexicale, syntaxique, sémantique, exécution) garantit la cohérence des messages d'erreur.

La hiérarchie d'exceptions bien structurée distingue clairement trois catégories d'erreurs :
\begin{itemize}
    \item \texttt{SyntaxException} pour les erreurs de parsing (lexicales et syntaxiques)
    \item \texttt{SemanticException} pour les erreurs d'analyse sémantique
    \item \texttt{JajaCodeRuntimeException} et ses sous-classes pour les erreurs d'exécution
\end{itemize}

L'intégration étroite entre \texttt{DiagnosticCollector} et les exceptions de compilation permet de fournir des messages d'erreur riches et détaillés, incluant tous les problèmes détectés plutôt que seulement le premier.
Les tests exhaustifs, incluant \texttt{ExceptionTypeTest} pour vérifier le type d'exception levé, assurent la fiabilité du système.

Cette approche permet aux utilisateurs de comprendre rapidement les problèmes dans leur code et de les corriger efficacement, tout en facilitant la maintenance et l'évolution future du système de gestion des erreurs.


% ============================================================================
% SECTION : Utilisation des outils de suivi de projet
% ============================================================================
\section{Utilisation des outils de suivi de projet}

\subsection{Introduction}
Cette section détaille l'utilisation des outils de suivi de projet mis en place pour le développement du compilateur MiniJAJA. Nous présentons ici l'intégration et l'utilisation de GitLab CI/CD, Nexus Repository Manager et SonarQube dans notre processus de développement.

\subsection{Vue d'ensemble de la chaîne CI/CD}

\subsubsection{Architecture globale}

Le projet MiniJAJA utilise une chaîne d'intégration continue complète basée sur GitLab CI/CD, intégrant Maven, JaCoCo, SonarQube et Nexus. Le pipeline automatise l'ensemble du processus, de la compilation jusqu'au déploiement des artefacts.

\begin{figure}[h]
    \centering
    \begin{verbatim}
    Git Push → Build → Tests → SonarQube → Package → Deploy (Nexus)
                 ↓       ↓        ↓
              compile  JUnit   Quality    JAR    Repository
                      JaCoCo   Gates
    \end{verbatim}
    \caption{Architecture du pipeline CI/CD}
\end{figure}

\subsubsection{Objectifs et bénéfices}

Les outils de CI/CD permettent de :
\begin{itemize}
    \item Automatiser les builds et tests à chaque commit
    \item Garantir la qualité du code via l'analyse statique SonarQube
    \item Centraliser les artefacts de build dans Nexus
    \item Détecter rapidement les régressions
    \item Faciliter les releases avec maven-release-plugin
    \item Assurer la reproductibilité des builds
\end{itemize}

% ============================================================================
% GITLAB CI/CD
% ============================================================================
\subsection{GitLab CI/CD}

\subsubsection{Configuration du pipeline}

\paragraph{Fichier .gitlab-ci.yml}

Le fichier \texttt{.gitlab-ci.yml} définit l'ensemble du pipeline. Il utilise l'image Docker \texttt{maven:3.9.9-eclipse-temurin-21-alpine} pour garantir un environnement de build cohérent.

Configuration globale :
\begin{lstlisting}[language=yaml, caption=Configuration globale GitLab CI]
default:
  image: maven:3.9.9-eclipse-temurin-21-alpine
  tags:
    - avm_groupe7
  before_script:
    - chmod +x settings.sh && ./settings.sh
  retry:
    max: 2
    when:
      - api_failure
      - script_failure
      - unknown_failure
\end{lstlisting}

Le script \texttt{settings.sh} génère automatiquement le fichier \texttt{.settings.xml} Maven avec les credentials pour Nexus et SonarQube à partir des variables d'environnement sécurisées.

\paragraph{Variables d'environnement}

\begin{lstlisting}[language=yaml, caption=Variables Maven]
variables:
  MAVEN_OPTS: "-Dmaven.repo.local=.m2/repository"
  MAVEN_CLI_OPTS_SINGLE_THREAD: "-B -V -ntp -s .settings.xml
                                 --fail-at-end"
  MAVEN_CLI_OPTS: "${MAVEN_CLI_OPTS_SINGLE_THREAD} -T 1C"
\end{lstlisting}

Options utilisées :
\begin{itemize}
    \item \texttt{-B} : mode batch (non-interactif)
    \item \texttt{-V} : affiche la version de Maven
    \item \texttt{-ntp} : pas de barre de progression (logs plus propres)
    \item \texttt{-T 1C} : compilation parallèle (1 thread par core CPU)
    \item \texttt{--fail-at-end} : continue le build malgré les échecs de modules
\end{itemize}

\paragraph{Cache Maven}

Le cache Maven améliore significativement les temps de build :
\begin{lstlisting}[language=yaml, caption=Configuration du cache]
cache:
  key: "m2-${CI_COMMIT_REF_SLUG}"
  policy: pull-push
  paths:
    - .m2/repository
  when: on_success
\end{lstlisting}

Le cache est :
\begin{itemize}
    \item \textbf{Partagé par branche} : clé basée sur \texttt{CI\_COMMIT\_REF\_SLUG}
    \item \textbf{Pull-push} : téléchargé au début, uploadé à la fin
    \item \textbf{Conditionnel} : seulement si le job réussit
\end{itemize}

\paragraph{Stages du pipeline}

Le pipeline est organisé en 7 stages séquentiels :
\begin{enumerate}
    \item \textbf{.pre} : vérification des variables d'environnement
    \item \textbf{build} : compilation du projet Maven
    \item \textbf{test} : exécution des tests unitaires et d'intégration
    \item \textbf{sonarqube-check} : analyse de la qualité du code
    \item \textbf{package} : création des artefacts JAR
    \item \textbf{deploy} : déploiement vers Nexus
    \item \textbf{release} : préparation et exécution des releases
\end{enumerate}

\subsubsection{Jobs et leur configuration}

\paragraph{Job check-envvars (.pre)}

Ce job préliminaire vérifie la présence de toutes les variables d'environnement nécessaires :
\begin{lstlisting}[language=yaml, caption=Verification des variables d'environnement]
check-envvars:
  stage: .pre
  script:
    - >
      if [[ -z "$MAVEN_ENT_PASSWORD" ]]
      || [[ -z "$MAVEN_ENT_USERNAME" ]]
      || [[ -z "$MAVEN_MASTER_PASSWORD" ]]
      || [[ -z "$SONAR_HOST_URL" ]]
      || [[ -z "$SONAR_TOKEN" ]];
      then echo "Missing env vars"; exit 1; fi
\end{lstlisting}

Variables vérifiées :
\begin{itemize}
    \item \texttt{MAVEN\_ENT\_USERNAME/PASSWORD} : authentification Nexus
    \item \texttt{MAVEN\_MASTER\_PASSWORD} : encryption Maven
    \item \texttt{SONAR\_HOST\_URL/TOKEN} : connexion SonarQube
\end{itemize}

\paragraph{Job build}

\begin{lstlisting}[language=yaml, caption=Job de build]
build:
  stage: build
  rules:
    - if: $CI_PIPELINE_SOURCE == 'merge_request_event'
    - if: $CI_COMMIT_TAG
    - if: $CI_COMMIT_BRANCH == 'main'
    - if: $CI_COMMIT_BRANCH == 'dev'
  script:
    - mvn ${MAVEN_CLI_OPTS} clean compile
\end{lstlisting}

Ce job :
\begin{itemize}
    \item Compile tous les modules Maven
    \item S'exécute sur les branches \texttt{main}, \texttt{dev}, les tags et les merge requests
    \item Utilise la compilation parallèle pour optimiser le temps
    \item Nettoie les artefacts précédents avec \texttt{clean}
\end{itemize}

\paragraph{Job test}

\begin{lstlisting}[language=yaml, caption=Job de test]
test:
  stage: test
  script:
    - mvn ${MAVEN_CLI_OPTS} -pl '!GUI,!CoverageReport'
          jacoco:prepare-agent test jacoco:report
  artifacts:
    reports:
      junit: "**/target/surefire-reports/TEST-*.xml"
      coverage_report:
        coverage_format: jacoco
        path: "**/target/site/jacoco/jacoco.xml"
    paths:
      - "**/target/surefire-reports/TEST-*.xml"
      - "**/target/site/jacoco/jacoco.xml"
      - "**/target/classes/"
    expire_in: 1 week
\end{lstlisting}

Ce job :
\begin{itemize}
    \item Exclut les modules \texttt{GUI} et \texttt{CoverageReport}
    \item Utilise JaCoCo pour mesurer la couverture de code
    \item Génère des rapports JUnit pour l'intégration GitLab
    \item Conserve les artefacts pendant 1 semaine
    \item Publie automatiquement les métriques de couverture dans l'interface GitLab
\end{itemize}

\paragraph{Job testIHM}

Test spécifique pour l'interface graphique avec environnement X11 virtuel :

\begin{lstlisting}[language=yaml, caption=Job de test IHM avec Xvfb]
testIHM:
  stage: test
  before_script:
    - chmod +x settings.sh && ./settings.sh
    - apk add xvfb libx11-dev xvfb-run gtk+3.0
    - export DISPLAY=:99
    - Xvfb $DISPLAY -screen 0 1400x900x24 +extension RANDR &
  script:
    - mvn ${MAVEN_CLI_OPTS} -pl 'LexerParser' -am install
          -DskipTests
    - xvfb-run -a mvn ${MAVEN_CLI_OPTS} -pl 'GUI'
          jacoco:prepare-agent test jacoco:report
\end{lstlisting}

Ce job :
\begin{itemize}
    \item Installe Xvfb (X Virtual FrameBuffer) pour les tests GUI headless
    \item Configure un écran virtuel 1400x900
    \item Compile d'abord les dépendances (\texttt{LexerParser})
    \item Exécute les tests GUI dans l'environnement virtuel
\end{itemize}

\paragraph{Job package}

\begin{lstlisting}[language=yaml, caption=Job de package]
package:
  stage: package
  script:
    - mvn ${MAVEN_CLI_OPTS} -DskipTests package
  artifacts:
    paths:
      - "**/target/*.jar"
    expire_in: 1 week
  needs:
    - test
    - testIHM
\end{lstlisting}

Ce job :
\begin{itemize}
    \item Génère les fichiers JAR de tous les modules
    \item Skip les tests (déjà exécutés dans le stage précédent)
    \item Dépend explicitement de \texttt{test} et \texttt{testIHM}
    \item Conserve les JARs pendant 1 semaine
\end{itemize}

\paragraph{Job sonarqube-check}

\begin{lstlisting}[language=yaml, caption=Job d'analyse SonarQube]
sonarqube-check:
  stage: sonarqube-check
  variables:
    SONAR_USER_HOME: "${CI_PROJECT_DIR}/.sonar"
    GIT_DEPTH: "0"
  cache:
    key: "sonar-cache"
    paths:
      - .sonar/cache
  script:
    - mvn ${MAVEN_CLI_OPTS_SINGLE_THREAD} -DskipTests
          sonar:sonar
          -Dsonar.projectKey=avm-2025-groupe-7
          -Dsonar.projectName='avm-2025-groupe-7'
  needs:
    - build
    - test
    - testIHM
\end{lstlisting}

Ce job :
\begin{itemize}
    \item Clone l'historique Git complet (\texttt{GIT\_DEPTH=0}) pour l'analyse
    \item Utilise un cache SonarQube séparé
    \item Exécute en single-thread pour éviter les problèmes de concurrence, car le plugin maven \texttt{sonar-maven-plugin} n'est pas thread-safe
    \item Attend les résultats des tests pour avoir les rapports de couverture
\end{itemize}

\paragraph{Job deploy}

\begin{lstlisting}[language=yaml, caption=Job de deploiement Nexus]
deploy:
  stage: deploy
  retry: 2
  rules:
    - if: $CI_COMMIT_TAG
    - if: $CI_COMMIT_BRANCH == 'main'
  needs:
    - package
  script:
    - mvn ${MAVEN_CLI_OPTS} -DskipTests deploy
\end{lstlisting}

Ce job :
\begin{itemize}
    \item S'exécute uniquement sur \texttt{main} et les tags
    \item Réessaye jusqu'à 2 fois en cas d'échec réseau
    \item Déploie les artefacts vers Nexus
\end{itemize}

\paragraph{Jobs release}

Le processus de release est divisé en deux jobs manuels :

\begin{lstlisting}[language=yaml, caption=Release prepare]
release:prepare:
  stage: release
  rules:
    - if: $CI_COMMIT_BRANCH == 'main'
      when: manual
  before_script:
    - apk add git openssh-client
    - git config --global user.email "ci@groupe7.fr"
    - git config --global user.name "CI Release Bot"
    # Configuration SSH pour push
  script:
    - mvn ${MAVEN_CLI_OPTS_SINGLE_THREAD}
          release:clean release:prepare
          -DpushChanges=true -Darguments="-DskipTests"
\end{lstlisting}

\begin{lstlisting}[language=yaml, caption=Release perform]
release:perform:
  stage: release
  rules:
    - if: $CI_COMMIT_BRANCH == 'main'
      when: manual
  needs:
    - release:prepare
  script:
    - mvn ${MAVEN_CLI_OPTS_SINGLE_THREAD}
          release:perform -Darguments="-DskipTests"
\end{lstlisting}

Ces jobs :
\begin{itemize}
    \item Sont déclenchés manuellement depuis GitLab
    \item Créent un tag Git avec la nouvelle version
    \item Mettent à jour les versions dans les POM
    \item Déploient la release vers Nexus
\end{itemize}

\subsubsection{Règles de déclenchement}

Les jobs utilisent des règles (\texttt{rules}) pour contrôler leur exécution :
\begin{itemize}
    \item \textbf{Merge Requests et branche dev} : tous les jobs sauf deploy et release
    \item \textbf{Branche main} : tous les jobs incluant deploy
    \item \textbf{Tags} : tous les jobs incluant deploy
\end{itemize}

\subsubsection{Infrastructure des runners GitLab}

L'infrastructure d'exécution du pipeline repose sur plusieurs runners GitLab, combinant ressources universitaires et infrastructure personnelle pour garantir disponibilité et performance.

\paragraph{Configuration standard}

La plupart des groupes de projet disposent d'un unique runner fourni par l'université. Ce runner est partagé entre plusieurs groupes, ce qui peut entraîner des temps d'attente lors des périodes de forte activité (avant les releases, lors de multiples merge requests simultanées).

\paragraph{Infrastructure renforcée du groupe}

Pour pallier ces limitations et garantir une disponibilité maximale du pipeline, le groupe a mis en place une infrastructure renforcée :

\begin{itemize}
    \item \textbf{1 runner universitaire} : fourni par l'université, partagé avec d'autres groupes
    \item \textbf{3 runners self-hosted} : hébergés sur un VPS (Virtual Private Server) personnel
\end{itemize}

Total : \textbf{4 runners} disponibles pour le projet.

\paragraph{VPS personnel}

Un membre du groupe loue pour ses activités personnelles un VPS chez Infomaniak. Ce serveur a été configuré pour héberger 3 runners GitLab CI/CD supplémentaires dédiés au projet. Cette infrastructure apporte plusieurs avantages :

\begin{itemize}
    \item \textbf{Disponibilité garantie} : les runners personnels prennent le relais si les runners universitaires sont indisponibles
    \item \textbf{Pas d'attente en période de rush} : avec 4 runners, les merge requests peuvent exécuter tous les jobs sans file d'attente, même lors des périodes critiques (avant une release, multiples MR simultanées)
    \item \textbf{Résilience} : redondance complète en cas de défaillance des infrastructures universitaires
\end{itemize}

\paragraph{Configuration des runners}

Les runners sont configurés avec les tags appropriés (ex: \texttt{avm\_groupe7}) pour être utilisés uniquement par le projet du groupe. Chaque runner :
\begin{itemize}
    \item Utilise Docker pour l'isolation des builds
    \item Possède un cache local Maven pour optimiser les temps de compilation
    \item Est configuré pour exécuter les mêmes jobs que les runners universitaires
\end{itemize}

\paragraph{Bénéfices constatés}

Cette infrastructure renforcée a permis de :
\begin{itemize}
    \item Garantir une très haute disponibilité du pipeline (continuité même en cas de maintenance ou de problème au sein du système universitaire)
    \item Permettre aux développeurs de pusher et tester leurs branches sans impact sur les autres membres de l'équipe
    \item Faciliter les releases avec des pipelines complets s'exécutant en moins de 5 minutes
\end{itemize}

% ============================================================================
% NEXUS
% ============================================================================
\subsection{Nexus Repository Manager}

\subsubsection{Présentation et rôle}

Nexus Repository Manager est utilisé comme gestionnaire central de dépôts Maven pour :
\begin{itemize}
    \item Stocker les artefacts de build (JAR, POM, sources)
    \item Gérer les versions SNAPSHOT et RELEASE
    \item Servir de proxy pour les dépendances externes (Maven Central)
    \item Faciliter la distribution des livrables entre équipes
    \item Assurer la reproductibilité des builds
\end{itemize}

URL Nexus fournis par l'université : \texttt{https://disc.univ-fcomte.fr/cr700-nexus/}

\subsubsection{Configuration des repositories}

\paragraph{Repository des releases}

Configuration dans le \texttt{pom.xml} parent :
\begin{lstlisting}[language=enhancedxml, caption=Configuration du repository releases]
<distributionManagement>
    <repository>
        <id>nexus-deptinfo</id>
        <url>https://disc.univ-fcomte.fr/cr700-nexus/
             repository/maven-releases/</url>
    </repository>
</distributionManagement>
\end{lstlisting}

Caractéristiques :
\begin{itemize}
    \item \textbf{Politique} : versions RELEASE uniquement (sans \texttt{-SNAPSHOT})
    \item \textbf{Déploiement} : immuable (une fois déployée, une version ne peut être écrasée)
    \item \textbf{Usage} : versions stables et livrables en production
\end{itemize}

\paragraph{Repository des snapshots}

\begin{lstlisting}[language=enhancedxml, caption=Configuration du repository snapshots]
<snapshotRepository>
    <id>nexus-deptinfo-snapshots</id>
    <url>https://disc.univ-fcomte.fr/cr700-nexus/
         repository/maven-snapshots/</url>
</snapshotRepository>
\end{lstlisting}

Caractéristiques :
\begin{itemize}
    \item \textbf{Politique} : versions SNAPSHOT uniquement (suffixe \texttt{-SNAPSHOT})
    \item \textbf{Déploiement} : mutable (peut être redéployé)
    \item \textbf{Usage} : développement continu, intégration continue
    \item \textbf{Rétention} : configurable (généralement limité dans le temps)
\end{itemize}

\paragraph{Authentification}

L'authentification vers Nexus est configurée dans le fichier \texttt{.settings.xml} généré dynamiquement par le script \texttt{settings.sh} :

\begin{lstlisting}[language=enhancedxml, caption=Authentification Nexus]
<servers>
    <server>
        <id>nexus-deptinfo</id>
        <username>${env.MAVEN_ENT_USERNAME}</username>
        <password>${env.MAVEN_ENT_PASSWORD}</password>
    </server>
    <server>
        <id>nexus-deptinfo-snapshots</id>
        <username>${env.MAVEN_ENT_USERNAME}</username>
        <password>${env.MAVEN_ENT_PASSWORD}</password>
    </server>
</servers>
\end{lstlisting}

Les credentials sont stockés en tant que variables d'environnement sécurisées dans GitLab CI/CD.

\subsubsection{Gestion des versions}

\paragraph{Workflow de versioning}

\begin{enumerate}
    \item \textbf{Développement} : version SNAPSHOT (ex: \texttt{1.1-SNAPSHOT})
    \item \textbf{Release prepare} : maven-release-plugin crée la version release (\texttt{1.1})
    \item \textbf{Tag Git} : création automatique du tag \texttt{v1.1}
    \item \textbf{Bump version} : passage à la prochaine version SNAPSHOT (\texttt{1.2-SNAPSHOT})
    \item \textbf{Release perform} : déploiement de la version \texttt{1.1} vers Nexus releases
\end{enumerate}

\paragraph{Configuration maven-release-plugin}

\begin{lstlisting}[language=enhancedxml, caption=Configuration du plugin release]
<plugin>
    <groupId>org.apache.maven.plugins</groupId>
    <artifactId>maven-release-plugin</artifactId>
    <version>3.1.1</version>
    <configuration>
        <tagNameFormat>v@{project.version}</tagNameFormat>
        <autoVersionSubmodules>true</autoVersionSubmodules>
        <releaseProfiles>release</releaseProfiles>
        <scmCommentPrefix>[Release] </scmCommentPrefix>
        <goals>deploy</goals>
    </configuration>
</plugin>
\end{lstlisting}

Options utilisées :
\begin{itemize}
    \item \texttt{autoVersionSubmodules} : synchronise les versions de tous les modules
    \item \texttt{tagNameFormat} : format du tag Git (\texttt{v1.1})
    \item \texttt{scmCommentPrefix} : préfixe pour les commits de release
\end{itemize}

\subsubsection{Avantages constatés}

\begin{itemize}
    \item \textbf{Centralisation} : tous les artefacts au même endroit
    \item \textbf{Traçabilité} : historique complet des versions déployées
    \item \textbf{Reproductibilité} : possibilité de rebuilder des versions anciennes
    \item \textbf{Performance} : cache local des dépendances externes
    \item \textbf{Sécurité} : contrôle d'accès aux artefacts
\end{itemize}

% ============================================================================
% SONARQUBE
% ============================================================================
\subsection{SonarQube}

\subsubsection{Présentation et objectifs}

SonarQube est utilisé pour l'analyse statique continue du code afin de :
\begin{itemize}
    \item Détecter les bugs potentiels et erreurs de code
    \item Identifier les failles de sécurité (vulnerabilities)
    \item Mesurer la dette technique (code smells)
    \item Évaluer la couverture de code par les tests
    \item Vérifier le respect des bonnes pratiques Java
    \item Détecter la duplication de code
\end{itemize}

URL SonarQube : \texttt{\$\{SONAR\_HOST\_URL\}} (configuré via variable d'environnement)

\subsubsection{Configuration du projet}

\paragraph{Propriétés SonarQube dans le POM}

Configuration dans le \texttt{pom.xml} parent :
\begin{lstlisting}[language=enhancedxml, caption=Configuration SonarQube]
<properties>
    <sonar.qualitygate.wait>true</sonar.qualitygate.wait>
    <sonar.coverage.jacoco.xmlReportPaths>
        ${project.basedir}/target/site/jacoco/jacoco.xml
    </sonar.coverage.jacoco.xmlReportPaths>
    <sonar.exclusions>
        LexerParser/src/main/java/.../JajaCodeInterpreter.java,
        LexerParser/src/main/java/.../MiniJajaCompiler.java
    </sonar.exclusions>
</properties>
\end{lstlisting}

Options importantes :
\begin{itemize}
    \item \texttt{qualitygate.wait=true} : le build échoue si la Quality Gate n'est pas passée
    \item \texttt{coverage.jacoco.xmlReportPaths} : chemin vers les rapports JaCoCo
    \item \texttt{exclusions} : fichiers principaux exclus de l'analyse (points d'entrée)
\end{itemize}

\paragraph{Plugin sonar-maven-plugin}

\begin{lstlisting}[language=enhancedxml, caption=Plugin SonarQube Maven]
<plugin>
    <groupId>org.sonarsource.scanner.maven</groupId>
    <artifactId>sonar-maven-plugin</artifactId>
    <version>5.2.0.4988</version>
</plugin>
\end{lstlisting}

\paragraph{Intégration JaCoCo}

JaCoCo génère les rapports de couverture pour SonarQube :
\begin{lstlisting}[language=enhancedxml, caption=Configuration JaCoCo]
<plugin>
    <groupId>org.jacoco</groupId>
    <artifactId>jacoco-maven-plugin</artifactId>
    <version>0.8.14</version>
    <executions>
        <execution>
            <goals>
                <goal>prepare-agent</goal>
            </goals>
        </execution>
        <execution>
            <id>report</id>
            <phase>verify</phase>
            <goals>
                <goal>report</goal>
            </goals>
        </execution>
        <execution>
            <id>report-aggregate</id>
            <phase>verify</phase>
            <goals>
                <goal>report-aggregate</goal>
            </goals>
        </execution>
    </executions>
</plugin>
\end{lstlisting}

JaCoCo :
\begin{itemize}
    \item Instrumente le bytecode Java pour mesurer la couverture
    \item Génère un rapport XML pour SonarQube
    \item Agrège les rapports de tous les modules
\end{itemize}

\paragraph{Quality Gates}

Les Quality Gates définissent les seuils de qualité minimum :
\begin{itemize}
    \item \textbf{Couverture} : pourcentage minimal de code testé
    \item \textbf{Duplication} : taux maximal de code dupliqué
    \item \textbf{Bugs} : nombre maximal de bugs tolérés
    \item \textbf{Vulnerabilities} : aucune vulnérabilité de sévérité élevée
    \item \textbf{Code Smells} : dette technique acceptable
\end{itemize}

Le pipeline échoue si une Quality Gate n'est pas satisfaite, empêchant le déploiement de code de mauvaise qualité.

\subsubsection{Métriques analysées}

\paragraph{Fiabilité (Reliability)}

Mesure la probabilité de défaillance du logiciel :
\begin{itemize}
    \item \textbf{Bugs} : erreurs de code susceptibles de causer des défaillances
    \item \textbf{Rating} : de A (excellent) à E (mauvais)
    \item Exemples détectés :
    \begin{itemize}
        \item NullPointerException potentielles
        \item Ressources non fermées (fichiers, connexions)
        \item Conditions toujours vraies/fausses
        \item Boucles infinies potentielles
    \end{itemize}
\end{itemize}

\paragraph{Sécurité (Security)}

Identifie les vulnérabilités et risques de sécurité :
\begin{itemize}
    \item \textbf{Vulnerabilities} : failles de sécurité confirmées
    \item \textbf{Security Hotspots} : code sensible nécessitant une revue manuelle
    \item \textbf{Rating} : de A (sécurisé) à E (très vulnérable)
    \item Exemples :
    \begin{itemize}
        \item Injection SQL potentielle
        \item Désérialisation non sécurisée
        \item Utilisation de cryptographie faible
        \item Exposition d'informations sensibles
    \end{itemize}
\end{itemize}

\paragraph{Maintenabilité (Maintainability)}

Évalue la facilité de maintenance du code :
\begin{itemize}
    \item \textbf{Code Smells} : problèmes de conception et mauvaises pratiques
    \item \textbf{Dette technique} : temps estimé pour corriger tous les smells
    \item \textbf{Rating} : de A (très maintenable) à E (difficilement maintenable)
    \item \textbf{Duplication} : pourcentage de code dupliqué
    \item Exemples de code smells :
    \begin{itemize}
        \item Méthodes trop longues ou trop complexes (complexité cyclomatique élevée)
        \item Classes avec trop de responsabilités
        \item Code mort (jamais exécuté)
        \item Commentaires TODO non résolus
        \item Magic numbers (constantes non nommées)
    \end{itemize}
\end{itemize}

\paragraph{Couverture (Coverage)}

Mesure du code testé par les tests unitaires :
\begin{itemize}
    \item \textbf{Couverture de lignes} : pourcentage de lignes exécutées
    \item \textbf{Couverture de branches} : pourcentage de branches conditionnelles testées
    \item \textbf{Nouveau code} : couverture spécifique aux modifications récentes
    \item Rapports générés par JaCoCo et agrégés par SonarQube
\end{itemize}

\subsubsection{Intégration dans le pipeline CI/CD}

\paragraph{Analyse automatique}

L'analyse SonarQube est déclenchée automatiquement :
\begin{itemize}
    \item À chaque push sur \texttt{dev}
    \item Sur toutes les merge requests
    \item Exécution via : \texttt{mvn sonar:sonar}
    \item Authentification via token stocké dans \texttt{SONAR\_TOKEN}
\end{itemize}

\paragraph{Feedback et rapports}

Après chaque analyse :
\begin{itemize}
    \item Les résultats sont consultables sur l'interface SonarQube
    \item Le statut de la Quality Gate est affiché dans GitLab
    \item Le pipeline échoue si la Quality Gate n'est pas passée
    \item Les développeurs sont notifiés des problèmes détectés
\end{itemize}

\subsubsection{Retour d'expérience}

\paragraph{Bénéfices constatés}

\begin{itemize}
    \item \textbf{Détection précoce} : bugs trouvés avant la mise en production
    \item \textbf{Qualité constante} : maintien d'un niveau de qualité élevé
    \item \textbf{Formation} : les développeurs apprennent des erreurs détectées
    \item \textbf{Documentation} : SonarQube explique pourquoi chaque problème est important
    \item \textbf{Réduction de la dette} : visualisation et suivi de la dette technique
\end{itemize}

\paragraph{Difficultés rencontrées}

\begin{itemize}
    \item \textbf{Faux positifs} : certaines règles nécessitent des ajustements
    \item \textbf{Temps d'analyse} : peut ralentir le pipeline (d'où l'utilisation du cache)
    \item \textbf{Exclusions} : nécessité d'exclure les fichiers générés (ANTLR)
    \item \textbf{Seuils} : définition des Quality Gates adaptées au projet
\end{itemize}

% ============================================================================
% CONCLUSION
% ============================================================================
\subsection{Conclusion et perspectives}

\subsubsection{Bilan global}

L'intégration de GitLab CI/CD, Nexus et SonarQube dans le projet MiniJAJA a apporté :
\begin{itemize}
    \item \textbf{Automatisation complète} : du commit au déploiement
    \item \textbf{Qualité garantie} : Quality Gates empêchent le déploiement de code défaillant
    \item \textbf{Traçabilité} : historique complet des builds et artefacts
    \item \textbf{Collaboration facilitée} : feedback immédiat sur les merge requests
    \item \textbf{Compétences acquises} : maîtrise des outils DevOps professionnels
\end{itemize}

\subsubsection{Métriques du pipeline}

Statistiques actuelles :
\begin{itemize}
    \item \textbf{Temps moyen de build complet} : 3-4 minutes (avec cache)
    \item \textbf{Parallélisation} : utilisation de tous les cores CPU (\texttt{-T 1C})
    \item \textbf{Cache hit rate} : \textasciitilde90\% (réduction significative du temps)
    \item \textbf{Retry policy} : jusqu'à 2 tentatives pour les échecs transitoires
    \item \textbf{Modules} : 7 modules Maven + 1 module documentation
\end{itemize}

\subsubsection{Améliorations possibles}

\begin{itemize}
    \item \textbf{Pipeline optimization} :
    \begin{itemize}
        \item Utilisation de stages DAG (Directed Acyclic Graph) pour plus de parallélisme
        \item Split des tests longs en jobs parallèles
        \item Cache Docker layers pour l'image Maven
    \end{itemize}
    \item \textbf{Tests supplémentaires} :
    \begin{itemize}
        \item Tests end-to-end automatisés
        \item Tests de performance (benchmarks)
        \item Tests de sécurité (OWASP Dependency Check)
    \end{itemize}
    \item \textbf{Métriques avancées} :
    \begin{itemize}
        \item Intégration avec des outils de monitoring
        \item Dashboards de métriques de qualité
        \item Alertes automatiques sur dégradation
    \end{itemize}
    \item \textbf{Déploiement continu} :
    \begin{itemize}
        \item Déploiement automatique vers environnement de staging
        \item Blue/green deployment
        \item Rollback automatique en cas d'échec
    \end{itemize}
\end{itemize}

\subsubsection{Recommandations}

Pour de futurs projets similaires :
\begin{itemize}
    \item \textbf{Commencer tôt} : mettre en place la CI/CD dès le début du projet
    \item \textbf{Quality Gates strictes} : ne pas tolérer la dégradation de qualité
    \item \textbf{Cache agressif} : optimiser les temps de build avec du cache
    \item \textbf{Tests en priorité} : écrire les tests avant de coder (TDD)
    \item \textbf{Documentation} : documenter la configuration du pipeline
    \item \textbf{Monitoring} : surveiller les temps de build et optimiser continuellement
    \item \textbf{Sécurité} : utiliser des variables d'environnement pour les secrets
\end{itemize}


\end{document}
